\begin{definition}[The implementation of a protocol]
  \label{def:aba-implementation}
  \lean{PLTS.ABA.Implementation.NetworkEvent}\lean{PLTS.ABA.Implementation.actsAt}\lean{PLTS.ABA.Implementation.coinLabelMap}\lean{PLTS.ABA.Implementation.RoundRecordMap}\lean{PLTS.ABA.Implementation.IsRoundRecord}\lean{PLTS.ABA.Implementation.NetworkState}\lean{PLTS.ABA.Implementation.roundOf}\lean{PLTS.ABA.Implementation.NetworkState.writeGhost}\lean{PLTS.ABA.Implementation.roundOwn}\lean{PLTS.ABA.Implementation.ProgramStep}\lean{PLTS.ABA.Implementation.NetworkStep}\lean{PLTS.ABA.Implementation.IsRoundRuleTable}\lean{PLTS.ABA.Implementation.system}\leanok
  \uses{def:aba-labels, def:aba-parameters, def:aba-round-loop-state, def:relabel, def:synchronised-product, def:process-algebra, def:idle-family}
  A protocol as it runs: $n$ programs, one per process, beside one network holding the message sets, the DECIDED sets and the corrupted set with its budget, beside the coin oracle.
  The system is parametric in the graded-agreement implementation --- its message type, its per-process per-round round record, and its rows --- and in a per-round ghost record the adversary holds, with an update \leandecl{PLTS.ABA.Implementation.NetworkState.writeGhost} applied on every row and an output relation on the announced bit guarding the two graded-agreement return rows (D30), so the round loop, the DECIDED sets, the coin handshake, corruption, the adversary and the composition pipeline are written once: \leandecl{PLTS.ABA.Implementation.system}.
  What an implementation supplies about its own rows, and what the lemmas that read a row off its label consume, is \leandecl{PLTS.ABA.Implementation.IsRoundRuleTable}.
\end{definition}
